\section{SUMMARY AND CONCLUSION}

This project developed the Digital Farmers' Produce System, a cloud-based platform designed to support crop management and market linkages in the Municipality of Sinacaban, Misamis Occidental. The project responded to practical problems identified among local farmers, including limited price information, difficulty finding consistent buyers, delayed transactions, and the absence of structured barangay-level crop visibility. These local problems reflect broader national concerns on post-harvest loss, information asymmetry, and uneven digital inclusion in agriculture.

Using the Rapid Application Development (RAD) framework, the project translated user needs into a role-based web application for farmers, buyers, Department of Agriculture/Municipal Agriculture Office (DA/MAO) staff, and administrators. Farmers can post available produce, manage listings, and request agricultural resources. Buyers can browse available products, search listings, and contact farmers. DA/MAO staff can monitor listings, manage requests, view summaries, and export data for reporting. The system was implemented using a PHP-MySQL stack, the Flight microframework, Bootstrap-based responsive interfaces, and shared cloud hosting. Security and reliability were addressed through HTTPS, role-based access control, password hashing, input validation, prepared statements, backups, and planned security testing.

Pilot user acceptance results were favorable. Across farmers, buyers, and DA/MAO officers, the mean ratings for perceived usefulness, perceived ease of use, attitude toward using, behavioral intention to use, system functionality, and overall satisfaction all fell within the \textit{Strongly Agree} range. Buyers produced the highest overall mean, followed by DA/MAO officers and farmers. The farmer results remained positive, but they also indicate the need for training, assistance, and sustained buyer participation to convert initial acceptance into regular system use.

The project demonstrates that a localized, user-centered agricultural platform can provide a practical bridge between farmers, buyers, and the local agriculture office. Its value lies not only in advertising produce but also in making agricultural activity visible, searchable, and usable for support planning. The system should therefore be viewed as a municipal coordination tool that can strengthen market linkage and decision support when paired with user training, reliable data entry, buyer engagement, and security validation.

In conclusion, the Digital Farmers' Produce System achieved its primary development and pilot-evaluation goals. It offers a feasible, low-cost, and context-sensitive approach to digital agricultural coordination in Sinacaban. While wider deployment and long-term impact assessment are still needed, the pilot results support the continued refinement and gradual implementation of the system as part of local agricultural modernization efforts.